\chapter{Advanced: runs, scenarios and aggregate planning}\label{ch:advanced}

The \ui{Advanced} group of the sidebar is for engineers and analysts: it shows every solver run, lets you run a
single algorithm on any data, and holds the research datasets used to compare algorithms. Day-to-day planning
does not need it.

\section{Runs}\label{sec:runs}

\shot{runs}{The Runs page: every solver run with status, scenario, algorithm, result and solve time.}{runs}

The list shows every run of every period and scenario: status, the data it ran on (with its short ID), the
algorithm and seed and budget, the solver result, the data revision, the solve time and when it was created.
Search by name, ID or input hash, and filter by status, algorithm or scenario. \ui{New single run} starts one
run (Section~\ref{sec:single-run}); \ui{Scenarios} opens the JSON scenarios. The list shows the most recent runs.

Click a run to open its details:

\shot{run_detail}{Details of the hybrid run: progress, objective, key figures and a validated Gantt chart.}{run_detail}

\begin{itemize}
  \item \textbf{Run progress}: queued, solver started (and how long it waited for a slot), finished and
        validated (and how long it solved). A running run refreshes every 1.5 seconds; it is safe to leave
        the page.
  \item \textbf{Figures}: objective score, required-lot completion, schedule end (including build-ahead),
        weighted tardiness, setup time, idle time, safety-stock shortfall and productive utilisation.
  \item \textbf{Gantt by machine}, with the number of operations that passed the independent validator.
  \item \textbf{Reproducibility}: the data source (period, scenario, aggregate plan or one-off JSON), the snapshot
        revision, the input hash, the solver status and time origin, and technical details.
  \item \ui{Open in scenario} goes to the period or scenario that owns the run; \ui{Retry with the same data}
        creates a new run from the same snapshot (for example after a failure). Hand edits cannot be retried.
\end{itemize}

\begin{note}
A schedule is only stored after it passes the independent validator, which checks every machine, shift,
setup, mold cycle, stock level and delivery again, regardless of which algorithm built the schedule.
\end{note}

\section{Creating a single run}\label{sec:single-run}

\shot{run_new}{Create a run: data source, solver settings and a summary.}{run_new}

\begin{enumerate}
  \item \textbf{Choose a data source}: \ui{Master data} (a factory dataset --- a period or a scenario; the run stores
        the selected revision), \ui{Aggregate plan} (the snapshot locked in a plan) or \ui{JSON file} (a
        one-off input up to 5~MB, not saved as master data).
  \item \textbf{Configure the solver}: algorithm, budget in seconds (up to the backend's limit) and random
        seed. The same input and seed make the run reproducible.
  \item Check the summary and click \ui{Start scheduling}.
\end{enumerate}

To generate several options to compare and commit, use the Generate dialog of a period or scenario instead.

\section{JSON scenarios}

\shot{scenarios}{JSON scenarios: whole datasets imported from files.}{scenarios}

A JSON scenario is a complete dataset (schema~4 or schema~5) imported from a file, used for experiments and
algorithm comparisons --- for example the benchmark instances in the project's \code{dataset/} folder. Import by
clicking \ui{Import JSON} or dropping a file on the page; every relation is validated before anything is
written. For day-to-day production scheduling use planning periods, which only need orders.

\shot{scenario}{A four-stage scenario: overview with hints from the interface's own analysis.}{scenario}

A scenario page has the same decision tabs as a period (\ui{Options}, \ui{Schedule \& compare}, \ui{Commit},
\ui{Analysis}, \ui{History}) plus \ui{Overview} and \ui{Data}:

\begin{itemize}
  \item \textbf{Overview} shows the scheduling state and a readiness check with hints (a mold near its cycle limit,
        products starting below safety stock, a stage needing a large share of shift time). Hints do not block
        scheduling.
  \item \textbf{Data} shows products and routes, machines, molds and calendar, orders and lots, and the settings and
        objective. Four-stage scenarios can be edited field by field (each change creates a new revision);
        five-stage scenarios are edited by replacing the JSON file (\ui{More} $\rightarrow$ \ui{Replace from
        JSON file\dots}).
  \item \textbf{More} also offers downloading the JSON of the current revision, a single run, and deleting the
        scenario. Runs and snapshots already created are kept.
\end{itemize}

Scenarios that belong to a planning period (they are created automatically when a period is saved) are not
listed here; a link to one redirects to its period.

\section{Aggregate planning (four-stage model)}

\shot{planning}{Aggregate planning: plans created and the form to create a new plan.}{planning}

Aggregate planning balances demand against capacity per day or per week before detailed scheduling --- the first
tier of the two-tier pipeline, available for four-stage scenarios. To create a plan:

\begin{enumerate}
  \item choose the dataset (the plan is locked to its current revision);
  \item choose the start and end date (up to the data's horizon) and the bucket size (\ui{Daily} or
        \ui{Weekly}; demand is grouped by due date);
  \item click \ui{Create plan}. The result is an immutable snapshot, ready for detailed scheduling.
\end{enumerate}

\shot{planning_detail}{An aggregate plan: feasibility, demand and inventory by product, capacity load by stage.}{planning_detail}

The plan page states whether the plan is \ui{Feasible at aggregate capacity} or \ui{Needs adjustment before
release}, and shows demand and inventory per product (demand, opening stock, quantity to produce, lot output,
projected closing stock, safety stock), the capacity load of each stage (required versus available time), and
demand versus planned output over time. \ui{Detailed scheduling} starts runs on the plan's snapshot.

\begin{note}
Aggregate capacity excludes sequence-dependent setups and mold maintenance; the validated detailed schedule
remains the execution authority.
\end{note}
